\section{Components}
\label{sec:components}

The code is seven Python files in \code{src/} and four scripts in \code{scripts/}
(Table~\ref{tab:files}). Every function and class has a docstring. The dependencies are pinned in
\code{requirements.txt}: PyTorch 2.15.0 (CPU build), sentence-transformers 6.3.0, chromadb 1.5.9,
pypdf 7.1.0, cryptography 50.0.2, gradio 6.32.0 and pytest 9.2.0. The OpenRouter call uses only
the standard library (\code{urllib}).

\begin{table}[H]
  \centering\small
  \begin{tabularx}{\textwidth}{@{}l X >{\raggedright\arraybackslash}p{4.6cm}@{}}
    \toprule
    File & Role & Key constants \\
    \midrule
    \code{src/ingest.py} & Reads the PDFs and text files, cuts them into chunks with stable ids. & \code{CHUNK_TOKENS} 512, \code{OVERLAP_TOKENS} 102 \\
    \code{src/index.py} & Embeds the chunks into Chroma (resumable), opens the index, reads every chunk back. & \code{EMBED_MODEL}, \code{COLLECTION} \\
    \code{src/bm25.py} & Okapi BM25 over the chunks, written from scratch with the standard library. & \code{k1} 1.5, \code{b} 0.75 \\
    \code{src/rag.py} & Retrieval, fusion, both gates, prompt, OpenRouter call, citation check, command line. & \code{TOP_K} 5, \code{CANDIDATES} 20, \code{SIM_THRESHOLD} 0.6, \code{RERANK_THRESHOLD} 0.5 \\
    \code{src/rerank.py} & Cross-encoder scoring: the gate score and reordering. & \code{GATE_RERANKER} \code{minilm} \\
    \code{src/app.py} & Gradio page: question, model, mode, threshold, answer, sources table. & \code{DISCLAIMER}, \code{MODES} \\
    \code{src/evaluate.py} & Retrieval comparison, threshold sweep, evaluation of five systems. & \code{THRESHOLDS}, \code{JUDGE_MAX_TOKENS} 4000 \\
    \midrule
    \code{scripts/fetch_corpus.py} & Downloads the corpus, from the manifest or by topic search. & 15 MB cap per PDF \\
    \code{scripts/make_dataset_links.py} & Writes \code{docs/dataset_links.md} from the manifest. & \\
    \code{scripts/check_eval_set.py} & Checks that the 20 evidence quotes are on their recorded page. & \\
    \code{scripts/run_evaluation.py} & Runs one day's share of the evaluation, with a lock and a log. & \code{STEPS} \\
    \bottomrule
  \end{tabularx}
  \caption{Files of the project and their role.}
  \label{tab:files}
\end{table}

\paragraph{Ingestion and index (\code{ingest.py}, \code{index.py}).}
The chunker receives the tokenizer as an argument. In production it is the tokenizer of BGE-M3,
so a window is exactly 512 tokens of the embedding model; the tests pass a fake tokenizer with one
token per word. The chunk id \code{source|page|start_token} is the same on every run for the same
settings, which is what makes the build resumable: the chunks still to embed are the set difference
between the corpus ids and the ids Chroma already holds, and stored ids that no longer match a chunk
are deleted. A PubMed chunk has no page; it is stored as $-1$ in Chroma and read back as none.

\paragraph{BM25 (\code{bm25.py}).}
An inverted index maps each term to the chunks that contain it and how often. Tokens are runs of
lowercase letters and digits. The score is Okapi BM25 with $k_1 = 1.5$ and $b = 0.75$, and an idf
with 1 inside the logarithm, $\ln\bigl(1 + (N - n + 0.5)/(n + 0.5)\bigr)$, which stays positive for
a term found in more than half the chunks (the Lucene variant). The code follows two
build-your-own-x tutorials and Robertson and Zaragoza (2009); its docstring credits them.

\paragraph{Retrieval and generation (\code{rag.py}).}
\code{rrf} gives a chunk at rank $r$ (counted from 1) a score of $1/(60 + r)$ per list and sorts by
the sum, after Cormack, Clarke and Buettcher (SIGIR 2009). \code{call_llm} sends one chat request
with \code{temperature} 0 and \code{max_tokens} 600. Its first statement is \code{check_free}, which
raises unless the model id ends in \code{:free} or is \code{openrouter/free}. \code{load_context}
loads the key, the index, BGE-M3, BM25 and the gate cross-encoder, and exits with one line when the
key or the index is missing.

\paragraph{Cross-encoder (\code{rerank.py}).}
\code{gate_score} applies a sigmoid to the cross-encoder's raw score (a logit), so 0.5 is the
model's own boundary between relevant and not relevant. It scores the passages one at a time and,
with \code{stop_at}, stops at the first one that reaches it. \code{rerank} sorts hits by score; it
is used only by the retrieval comparison. Three models are registered: \code{minilm}
(\code{cross-encoder/ms-marco-MiniLM-L6-v2}, the gate), \code{medcpt}
(\code{ncbi/MedCPT-Cross-Encoder}) and \code{bge} (\code{BAAI/bge-reranker-v2-m3}).

\paragraph{App (\code{app.py}).}
A Gradio page on \url{http://127.0.0.1:7860} with a question box, a model menu limited to the five
free models of \code{FREE_MODELS}, a choice of mode (default: hybrid + cross-encoder gate), a
threshold slider from 0 to 1 that applies to the gate of the chosen mode, the answer with the
disclaimer, and a table of the five passages with their page, similarity and whether they were
cited. Changing the mode resets the slider to that gate's default (0.6 or 0.5). An empty question
and a failed model call become a message instead of a crash.

\paragraph{Evaluation (\code{evaluate.py}, \code{run_evaluation.py}).}
\code{run} is resumable: each answer is appended to \nolinkurl{results/records.jsonl} with the threshold
it was made with, and a rerun skips what is there. It refuses to mix records made at another
threshold. \code{run_evaluation.py} asks OpenRouter how many free requests are left today (the
question itself uses none), runs the five systems with the default model and then
\code{rag_rerank} with the second free model, each with \code{--judge}, and appends its output to
\code{results/logs/<UTC date>.log}.
